% Plano SCX no esqueleto ABNT (abntex2).
% Diferenças da v1: sem \imprimircapa, sem \imprimirfolhaderosto, sem \tableofcontents;
% bloco de título manual na primeira página.
% Compilar: pdflatex plano-abnt.tex; biber plano-abnt; pdflatex plano-abnt.tex; pdflatex plano-abnt.tex
\documentclass[12pt,oneside,a4paper,brazil,sumario=tradicional]{abntex2}
\usepackage{amsmath,amssymb}
\usepackage{indentfirst}
\usepackage{booktabs,array,longtable}
\usepackage{graphicx}
\usepackage[utf8]{inputenc}
\usepackage[T1]{fontenc}
\usepackage{csquotes}
\usepackage[style=abnt,backend=biber,backref=false,accite]{biblatex}
\ExecuteBibliographyOptions{uniquename=false}
\addbibresource{refs.bib}
% abntex2 já carrega babel, hyperref, geometry e setspace
% numeração de seções independente de capítulos (documento curto, sem \chapter)
\renewcommand{\thesection}{\arabic{section}}
\renewcommand{\thesubsection}{\thesection.\arabic{subsection}}

% METADADOS (informativos; sem capa e sem folha de rosto)
\instituicao{Universidade de São Paulo}
\autor{Danilo Lindemberg Amorim}
\titulo{Emergência de Redes Livres de Escala com Agentes Semânticos baseados em Modelos de Decisão Sistema 1}
\local{São Paulo, SP}
\data{2026}
\tipotrabalho{Plano de pesquisa (dissertação de mestrado)}
\preambulo{Plano de pesquisa apresentado ao Programa de Pós-Graduação em Modelagem em Sistemas Complexos da Universidade de São Paulo como requisito parcial de qualificação.}

\begin{document}
\selectlanguage{brazil}
\frenchspacing
\textual
\markboth{}{}

% BLOCO DE TÍTULO
\begin{center}
{\Large\bfseries Emergência de Redes Livres de Escala com Agentes Semânticos baseados em Modelos de Decisão Sistema 1}\\[2ex]
{\large Danilo Lindemberg Amorim}\\[1ex]
{\normalsize Plano de Pesquisa\\
Programa de Pós-Graduação em Modelagem em Sistemas Complexos}
\end{center}

% 1. INTRODUÇÃO
\section{Introdução e Justificativa}

O mecanismo clássico de explicação da emergência de redes livres de escala é a ligação preferencial de \textcite{barabasi1999}, na qual o crescimento contínuo da rede combinado à probabilidade de conexão proporcional ao grau dos nós já existentes gera, de forma estocástica, uma topologia cuja distribuição de graus segue a Lei da Potência. Já foi observado, no entanto, que o comportamento livre de escala pode emergir de mecanismos formativos alternativos \parencite{caldarelli2002, vasquez2003, ferrericancho2003, dsouza2007, papadopoulos2012}.

A necessidade de entendimento dos mecanismos formativos subjacentes a este fenômeno motivou o desenvolvimento de modelos de simulação, que usualmente empregam a Modelagem Baseada em Agentes (ABM) \parencite{epstein1999, namatame2016, skyrms2000, zimmermann2005}. Um recente desenvolvimento, a Modelagem Baseada em Agentes Generativos (GABM) \parencite{ghaffarzadegan2024}, passou a utilizar Modelos Grandes de Linguagem (LLMs) como motor de decisão dos agentes na formação de redes complexas \parencite{demarzo2023, piao2025, papachristou2025, ji2025, ji2025b}. Entre as justificativas apresentadas por adotantes da abordagem estão tornar os agentes mais críveis \parencite{park2023, adornetto2025}, dotá-los de senso comum e raciocínio contextual \parencite{vezhnevets2023}, aproximá-los da racionalidade limitada humana, dotada de vieses e heurísticas, e flexibilizar regras de comportamento que de outra forma ficariam sujeitas ao viés do modelador \parencite{lu2024, chopra2024, larooij2025}.

Todavia, os LLMs apresentam limitações que dificultam a simulação em larga escala e a validade externa dos experimentos neles baseados. Sob as configurações usuais de decodificação, cada inferência é estocástica devido à amostragem de \textit{tokens} de saída \parencite{zhao2023,minaee2024}, além de apresentarem alto custo financeiro e latência de chamada \parencite{chopra2024,taillandier2026}. Como resultado, muitos estudos dependem de uma única execução \parencite{larooij2025}. Porém uma classe de modelos de IA pode suprir esta deficiência, os chamados Modelos de Decisão de Sistema 1 (DS1). O rótulo "Sistema 1" decorre da analogia com a decisão intuitiva descrita por \textcite{kahneman2003}, posteriormente formalizada na literatura de IA \parencite{bengio2019, booch2021, goyal2022}. O termo ganhou difusão comercial com o anúncio do modelo Jev \parencite{typesafe2026}, unificando um vocabulário antes disperso, sendo adotado pelos principais fornecedores \parencite{microsoft2026, cloudflare2026, openai2026}, já apresentando recente reflexo na academia \parencite{santos2026,du2026reliable,azizi2026labelsoverridedefinitionsjevstyle}.

Em vista disso, este trabalho propõe uma abordagem de Modelagem Baseada em Agentes Semânticos (SABM) se apoiando em modelos DS1 como motor de decisão dos agentes.  Investiga-se se essa abordagem é capaz de reproduzir as propriedades macroscópicas de redes livres de escala a partir de decisões microscópicas autônomas, e se seus resultados apresentam verossimilhança quando comparados aos modelos clássicos de formação de redes. O estudo se justifica por aplicar uma classe de modelos ainda pouco explorada em Sistemas Complexos e em ABM. A hipótese é de que o SABM atenda a grande parte das demandas que motivaram a GABM, porém, com custos computacionais, financeiros e temporais muito menores, além de habilitar a reprodutibilidade externa dos estudos quando o motor de decisão tiver pesos abertos. O laboratório de análise será a distribuição de popularidade (mensurada em "estrelas") de repositórios de projetos da rede sociotécnica do GitHub, que apresenta comportamento livre de escala.

% 2. OBJETIVOS
\section{Objetivos}

O propósito deste trabalho é validar empiricamente a metodologia SABM como ferramenta para simular a emergência de redes livres de escala, comparando seus resultados com o mecanismo clássico de \textcite{barabasi1999} e identificando benefícios, limitações e oportunidades de aprimoramento. Os objetivos específicos são:

\begin{itemize}
  \item Implementar e executar simulações de SABM reproduzindo o cenário de formação de redes de repositórios, em que a popularidade emerge de decisões microscópicas autônomas dos agentes desenvolvedores, os usuários da plataforma.
  \item Comparar quantitativamente as topologias simuladas com a topologia empírica, com o modelo Barabási-Albert, em suas versões unipartida e bipartida, e com mecanismos formativos alternativos, por meio da distribuição de grau dos repositórios e de medidas de agrupamento próprias de redes bipartidas.
  \item Investigar se a SABM reproduz a ligação preferencial linear ou se produz mecanismos emergentes distintos, como homofilia e otimização, que também convergem para a Lei da Potência, separando o efeito da exposição à popularidade do efeito da decisão semântica.
  \item Avaliar criticamente os benefícios e custos da SABM em relação à ABM, à GABM e aos modelos clássicos, considerando escalabilidade computacional, interpretabilidade, determinismo, reprodutibilidade e transparência das decisões.
  \item Propor refinamentos e diretrizes metodológicas para o uso da SABM em simulações de formação de redes complexas, sustentados pelas evidências empíricas obtidas.
\end{itemize}

% 3. REFERENCIAL
\section{Revisão Bibliográfica}

\subsection{Redes complexas e mecanismos de formação}

Esta seção sintetiza os modelos de redes complexas e os mecanismos propostos para explicar a formação de redes livres de escala, com base nas principais referências da ciência de redes \parencite{barabasi2002, newman2010, latora2017}.

Uma rede é representável por um grafo, com nós interligados por arestas que definem interações \parencite{newman2003, wang2003}. É dita complexa quando o comportamento coletivo não é previsível analisando apenas os componentes individuais, configurando emergência \parencite{mata2020}.

O ponto de partida do estudo das Redes Complexas é o modelo aleatório de \textcite{erdos1959}, que substituiu o paradigma do estudo de grafos determinísticos vigente até então, pela análise estatística de espaços de grafos gerados por sorteios de arestas \parencite{barabasi2002, latora2017}. Dadas conexões independentes e uniformes, os graus concentram-se próximos da média, e a formação do componente gigante exibe transição de fase no grau médio \parencite{albert2002}. Esses atributos fazem do modelo aleatório o referencial de mecanismos de formação, porque sem qualquer força motriz, não há razão para esperar qualquer heterogeneidade na rede \parencite{newman2003}.

O modelo de pequeno mundo de \textcite{watts1998} parte de uma rede regular, na qual cada nó se liga aos vizinhos mais próximos, e reconecta cada aresta com probabilidade $p$. Para $p = 0$, a rede permanece regular; para $p = 1$, torna-se aleatória. Em valores intermediários de $p$, a rede exibe caminhos médios curtos, típicos de grafos aleatórios, e alto agrupamento local, típico de redes regulares, com distribuição de grau ainda homogênea \parencite{latora2017}. Assim, o modelo delimita o regime no qual a conectividade eficiente coexiste com distribuição de grau igualitária.

Já o modelo livre de escala supõe distribuição de grau seguindo a Lei da Potência, a única função de densidade normalizável sob invariância a reescalonamento \parencite{albert2002, newman2005}, onde, $P(k) \sim k^{-\gamma}$.
Nesse modelo, a ausência de escala típica e a assimetria de graus impõem a formação de \textit{hubs}, nós com número extraordinário de conexões. Além disso, essas redes são descritas como robustas a falhas acidentais, porém frágeis diante de ataques direcionados aos \textit{hubs} \parencite{newman2003, wang2003}. O mecanismo formativo clássico é a ligação preferencial linear de \textcite{barabasi1999}, na qual, em crescimento contínuo, cada novo nó cria $m$ arestas cujo destino, o nó $i$, é sorteado pela probabilidade $\Pi(k_{i}) = k_{i} / \sum_{j} k_{j}$.

A combinação de crescimento e preferência proporcional ao grau reproduz distribuições com expoente $\gamma = 3$. A lei de crescimento do grau no tempo, $k_{i}(t) = m \left(\frac{t}{t_{i}}\right)^{1/2}$, formaliza a Vantagem do Pioneiro, pela qual o nó antigo acumula conexões de forma inevitável e nenhum recém-chegado o ultrapassa, independentemente de mérito \parencite{barabasi1999b}. As limitações do modelo são reconhecidas. A linearidade é um ponto de rigidez, pois no regime sublinear o modelo suprime os \textit{hubs} e cai em decaimento exponencial, e no superlinear colapsa em monopólio, em que um único nó absorve quase todas as arestas novas \parencite{albert2002}. Na forma generalizada $\Pi(k) \propto k^{\alpha}$, o regime linear corresponde a $\alpha = 1$, e o expoente $\alpha$ pode ser medido em redes reais pela taxa com que os nós adquirem arestas em função do grau \parencite{newman2001, jeong2003}. Além disso, o coeficiente de agrupamento do modelo tende a zero com o crescimento da rede, em dissonância com redes reais, nas quais ele permanece alto \parencite{newman2003}.

A linhagem de correções desse modelo é conhecida. \textcite{price1976}, em trabalho predecessor, formalizou em redes de citações, dirigidas e acíclicas, a vantagem cumulativa do efeito Mateus, já conhecida na Sociologia \parencite{merton1968}. O modelo Barabási-Albert é, em si, um caso especial não dirigido dessa formulação \parencite{newman2003}. Posteriormente, \textcite{bianconi2001} introduzem a aptidão individual fixada no nascimento do nó, tornando a probabilidade de conexão proporcional ao produto $\eta_{i}k_{i}$. Dele emergem o multi-escalonamento e a Vantagem do mais Apto, sob a qual recém-chegados aptos superam antigos, embora a topologia final fique presa às hipóteses da distribuição de aptidão \parencite{bianconi2001}. Este modelo também é criticado porque o cálculo de probabilidades pressupõe normalização global sobre o estado completo da rede, informação que um agente novo não possui \parencite{xu2010}. Paralelamente, \textcite{holme2002} acrescentam ao Barabási-Albert a formação controlada de tríades, estabilizando o coeficiente de agrupamento em valor finito ajustável enquanto preservam o expoente, também com o respaldo empírico em redes de colaboração científica \parencite{newman2003}.

Uma vertente alternativa, no entanto, deixa os modelos estocásticos e se concentra em abordagens de otimização. \textcite{dsouza2007} demonstram que a preferencialidade pode emergir de minimização local de custos. No modelo de pedágio na fronteira, cada novo nó equilibra o custo dos saltos intermediários na rede com a profundidade topológica até a raiz, e a saturação dos nós mais centrais gera rede livre de escala com expoente $\gamma = 2$ e cauda exponencial, por isso rotulada temperada. \textcite{ferrericancho2003} partem da otimização global estática, equilibrando eficiência de comunicação e custo de conexões; entre o extremo de ligação aleatória e o extremo de ligação forçada, o ótimo do processo é uma rede truncada de Lei da Potência, com o grafo em estrela como exceção teórica e inviável. Tais resultados abalam a tese formativa até aqui exclusivamente estocástica, pois a mesma assinatura topológica pode derivar de critérios deliberados, sem crescimento de nós.

O debate sobre a adequação do rótulo “livre de escala” segue vivo. \textcite{newman2003} admite desvios de cauda e defende o uso aproximativo do termo na prática. Já \textcite{broido2019}, ao analisar quase mil redes, verificam que distribuições alternativas, em particular a log-normal, retratam a maioria delas tão bem ou melhor do que a Lei da Potência. Todavia, propostas conciliadoras recentes mantêm o interesse em leis de cauda pesada \parencite{voitalov2019, holme2019}.

\subsection{Modelagem Baseada em Agentes}

A ABM representa o fenômeno por agentes individuais e suas interações \parencite{wilensky2015}, e indaga como interações locais de agentes autônomos e heterogêneos geram regularidade macroscópica \parencite{epstein1999}. Em contraste com a formulação de campo médio, que presume homogeneidade e substitui interações locais por interações uniformes de longo alcance, a ABM opera "de baixo para cima", lidando diretamente com indivíduos e com o caráter local das interações \parencite{boccara2004}.

Embora tenha sido inicialmente criticada pela suposta falta de rigor teórico \parencite{boccara2004}, adotantes dessa abordagem levantam argumentos a favor da sua robustez. \textcite{epstein1999} demonstra que a simulação equivale formalmente a um ato de dedução lógico, pela tese de Church-Turing, enquanto \textcite{namatame2016} se apoiam no princípio de irredutibilidade computacional de \textcite{wolfram2002} para defender que sistemas sociais complexos frequentemente apresentam comportamento computacionalmente irredutível, acessível apenas pela execução passo a passo. No entanto, reproduzir uma regularidade não prova as regras geradoras, como
advertido pelo princípio de equifinalidade de \textcite{bertalanffy1977}. É por isso que, nesse sentido, o modelador é aconselhado por \textcite{wilensky2015} a submeter o modelo a validação rigorosa e testes de sensibilidade para distinguir padrões recorrentes de artefatos de calibragem.

As principais limitações da ABM concentram-se nas regras de decisão dos agentes. Regras condicionais rígidas do tipo "se-então", arbitrárias, formuladas com base na intuição pessoal do modelador, muito sensíveis a alteração, ou pouco adaptáveis, estão entre as principais críticas apontadas à ABM \parencite{crooks2011, adornetto2025, chopra2024, gao2023, larooij2025, schulze2017}. Soma-se a isso a dificuldade de codificar processos mentais, mantendo a parcimônia do modelo \parencite{namatame2016}.

Nesse cenário, a GABM desponta como uma tentativa de preencher essa lacuna. Ela busca combinar a ABM tradicional com a IA Generativa ao se utilizar de um LLM como cérebro do agente. A premissa assumida é de que os LLMs possuem um conhecimento de mundo implícito e senso comum herdados do pré-treinamento em vastos corpos de dados humanos, o que permitiria substituir regras manuais rígidas por padrões comportamentais contextuais sociologicamente plausíveis, sem a necessidade de codificar cada reação individualmente \parencite{xi2023, wang2023, adornetto2025, ghaffarzadegan2024}. LLMs são modelos de redes neurais de arquitetura \textit{Transformer} escalonados a bilhões de parâmetros \parencite{vaswani2017, zhao2023}, cuja interação por \textit{prompts} dispensa ajuste fino para tarefas específicas sendo dotados de capacidades como aprendizado no contexto, cumprimento de instruções e raciocínio em múltiplas etapas \parencite{wei2022, zhao2023}.
Todavia, trabalhos publicados sob esta abordagem ainda são duramente criticados pela comunidade científica em questões que variam de rigor metodológico até limitações inerentes aos LLMs. Muitos dos estudos dessa classe se apoiam em conclusões sobre a validade de face da simulação fazendo juízo subjetivo da credibilidade do comportamento do agente, sem evidenciar a validade empírica do mecanismo simulado. Motivos apontados para este cenário são a opacidade da "caixa-preta" do LLM, que dificulta interpretar, padronizar e replicar os resultados; a estocasticidade das saídas, que fragiliza o teste reprodutível; e os vieses sociais e de seleção herdados do \textit{corpus} de treinamento, que caricaturam grupos humanos e induzem o modelo a reproduzir fenômenos já documentados na literatura, em vez de gerar a dinâmica que o modelador pretendia \parencite{larooij2025,taillandier2026}. Não obstante, custo e latência elevados por decisão também se colocam como obstáculos à reprodutibilidade dos estudos \parencite{chen2023}. Nesse sentido, a seguir é apresentada uma classe de modelos que tem potencial para mitigar algumas das limitações da GABM.

\subsection{Modelos de Decisão Sistema 1}

A expressão "Sistema 1" tem raízes no trabalho de \textcite{kahneman2003}, que descreveu dois modos de funcionamento cognitivo: o Sistema 1, intuitivo, rápido e automático, e o Sistema 2, deliberado e lento. Essa dualidade foi adotada na Inteligência Artificial como analogia para arquiteturas cognitivas \parencite{bengio2019, booch2021, goyal2022}. De modo geral, a geração autoregressiva de \textit{tokens} dos LLMs tende a aproximá-los do Sistema 1, ao passo que arquiteturas que acrescentam busca deliberada, planejamento ou módulos de verificação lhes conferem funções associadas ao Sistema 2 \parencite{booch2021, kambhampati2024, yao2023, sumers2024}. A analogia, contudo, não é exclusiva dos LLMs, pois a proposta de \textcite{bengio2019} tratava do aprendizado profundo em geral.

O custo e a latência da inferência dos LLMs motivaram arquiteturas em cascata, nas quais modelos mais baratos resolvem as consultas primeiro e somente as mais difíceis são encaminhadas a modelos maiores \parencite{chen2023}. Variações desse princípio incluem o roteamento entre um modelo forte e um fraco, treinado com dados de preferência \parencite{ong2025}, e a avaliação probabilística tipificada com escalonamento seletivo \parencite{nandakishore2026}. Nessas arquiteturas, a aceitação de uma resposta pode caber a um juiz, que avalia a saída de outro modelo, ou a um verificador, que a checa antes do uso; em FrugalGPT, por exemplo, uma função de avaliação decide se a resposta é aceita ou escalonada \parencite{chen2023}. Nos papéis de roteador e de triagem, a entrada semântica é interpretada e a decisão é tomada de imediato, sem debate deliberado.

Nesse contexto, os Modelos de Decisão Sistema 1 (DS1) constituem uma classe de modelos treinados para decisões tipificadas, com proabilidades calibradas. Cada requisição traz um estado em texto e perguntas cujo conjunto fechado de respostas é declarado de antemão, e o modelo devolve, numa única inferência, uma probabilidade para cada opção, sem gerar texto e, portanto, sem amostragem de \textit{tokens} \parencite{typesafe2026, laya2026}. O Jev, da TypeSafe AI \parencite{typesafe2026}, é o exemplo mais marcante desta classe, acompanhado de lançamentos recentes por outros fornecedores \parencite{microsoft2026, cloudflare2026, openai2026}.

A classe é heterogênea quanto à arquitetura e ao porte, como se observa nos modelos de pesos abertos Laya \parencite{laya2026} e Jebadiah \parencite{frontierinfra2026}. Os DS1 compartilham o formato de entrada e saída, com pergunta tipificada, opções fechadas e probabilidade por opção, que serve de critério para delimitar a classe. O juiz e o verificador descritos acima ficam fora da classe, pois o DS1 é o elemento de decisão inicial, cuja saída é aproveitável diretamente como decisão em um fluxo automatizado.

Suas propriedades pedem cautela por três razões. Primeiro, a calibração é um objetivo de treinamento, e não uma garantia: a documentação do Laya reporta erro de calibração esperado de 0,466 antes e de 0,081 após ajuste por temperatura, conforme relatado por \textcite{santos2026}. Segundo, a leitura da probabilidade de rótulos é sensível ao viés de rótulo e de ordem das opções em LLMs usados como classificadores \parencite{zhao2021}, e há relato de que, em modelos no estilo do Jev, os rótulos das opções podem prevalecer sobre suas definições \parencite{azizi2026labelsoverridedefinitionsjevstyle}. Terceiro, quando derivados de LLMs pré-treinados, os modelos DS1 podem herdar os vieses de corpus atribuídos à GABM. Soma-se a isso o fato de a classe ser de formação recente e documentada sobretudo por fornecedores e por \textit{preprints} \parencite{du2026reliable, santos2026, azizi2026labelsoverridedefinitionsjevstyle}, de modo que suas propriedades são tratadas aqui como hipóteses a verificar, e não como premissas.

Em síntese, frente ao LLM generativo, o DS1 abre mão da riqueza discursiva e da explicação em linguagem natural em troca de saída tipificada, probabilidade por opção e menor custo e latência por decisão \parencite{typesafe2026, laya2026}. Isso mitiga duas das limitações da GABM apontadas na seção anterior, a estocasticidade da amostragem e o custo por chamada, sem eliminar os vieses herdados do pré-treinamento nem garantir, por si só, determinismo numérico entre ambientes de execução.

\subsection{Plataformas de desenvolvimento como redes sociotécnicas}

O GitHub é uma das maiores plataformas de hospedagem e compartilhamento de código-fonte aberto do mundo \parencite{github2025, stackoverflow2025}, na qual desenvolvedores compartilham projetos, conhecidos como repositórios, e colaboram de maneira organizada e segura. Para expressar apreciação por um repositório, o desenvolvedor pode atribuir a ele uma estrela, um mecanismo análogo à curtida em redes sociais, de modo que a quantidade de estrelas é uma proxy de popularidade consolidada na literatura \parencite{borges2016, borges2018, gu2026}.

Trabalhos empíricos sobre essas plataformas relatam que a distribuição de estrelas por repositório segue a Lei da Potência, com uma elite de repositórios concentrando dezenas de milhares de estrelas, enquanto a maioria permanece próxima de zero \parencite{borges2018, gu2026}. Essa desigualdade é típica do efeito Mateus \parencite{merton1968}. A popularidade se forma no tempo, a partir de decisões repetidas de muitos desenvolvedores expostos a listas de descoberta com contagem visível, processo cumulativo análogo ao que a hipótese da ligação preferencial descreve \parencite{barabasi1999}.

Argumenta-se, com \textcite{orlikowski2008}, que nessas plataformas a tecnologia e a prática social são mutuamente constitutivas, o que as caracteriza como redes sociotécnicas. No lado social do acoplamento estão comunidades de uso, reputação, marketing técnico e círculos de divulgação; no lado técnico estão a adoção de linguagens e famílias tecnológicas, a infraestrutura de busca e ordenação, o README como superfície semântica de apresentação e os mecanismos de reutilização, como os forks \parencite{jiang2017}. O acoplamento também aparece em incidentes, como comprometimentos da cadeia de suprimento de pacotes de software \parencite{saayman2026, dholakia2026}, e sistemas de troca de estrelas, que fabricam ganho de popularidade, documentam que a métrica pode ser manipulada socialmente, embora seu registro seja técnico \parencite{he2025b}.

% 4. METODOLOGIA
\section{Metodologia}

A metodologia contempla simulação computacional com motor de decisão DS1, calibrada por dados empíricos do GitHub. O objeto observacional é um grafo bipartido, com desenvolvedores e repositórios públicos nas partições e arestas dirigidas de estrela emitidas ao longo do tempo; a popularidade de cada repositório é o seu grau de entrada no grafo bipartido, isto é, o número de estrelas recebidas. A definição do grafo orienta as comparações seguintes, pois no Barabási-Albert as arestas nascem com o nó novo, enquanto na rede de estrelas as duas partições crescem e as arestas ligam nós já existentes. Por isso, a comparação correta é feita com versões bipartidas dos modelos de crescimento \parencite{guillaume2006}, e o agrupamento é medido por coeficientes próprios de redes bipartidas \parencite{latapy2008, opsahl2013}, já que a projeção unipartida sobre os repositórios gera, por construção, uma clique para cada desenvolvedor e inflaria o agrupamento. Os dados observados servem de referencial de verossimilhança: se a emergência simulada reproduz o formato estatístico do fenômeno observado, a hipótese dos agentes semânticos ganha suporte; se revela mecanismos formativos distintos, o resultado enriquece o comparativo entre os mecanismos da seção 3.1.

O recorte de dados apoia-se em repositórios públicos criados entre janeiro de 2024 e abril de 2026, janela que acompanha cada repositório desde a criação. O GH Archive registra, de forma pública, granular e longitudinal, cada evento de estrela \parencite{grigorik_gharchive}, o que permite verificar o conjunto de dados em recortes temporais bem delimitados e confrontar a simulação com a topologia real observada. Por abranger o regime novo de contribuição assistida por LLMs \parencite{gu2026, popescu2026, taillandier2026}, a janela está sujeita a não estacionariedade e a estrelas fabricadas \parencite{he2025b}, monitoradas como ameaças à validade. São excluídos os repositórios vazios, os criados por bots, os nascidos de forks e os não-software, com a classificação por tipo conduzida pelo próprio motor DS1 sobre o conteúdo do README, com precisão e recall relatados em amostra de validação \parencite{jiang2017, he2025b}. A atividade mínima de dois commits nos primeiros 30 dias é um limite de corte a ser testado em análise de sensibilidade, e os repositórios arquivados permanecem na análise, já que continuam a receber estrelas. Para mitigar o viés de sobrevivência, extrai-se do histórico de versões do arquivo o README vigente em cada passo mensal; os repositórios apagados, presentes no GH Archive mas sem README recuperável, são contabilizados e sua fração é relatada. Os nomes e menções de repositórios são anonimizados, controlando a possibilidade de conhecimento prévio dos modelos sobre os repositórios, que podem integrar seus dados de treinamento; como a anonimização não impede a memorização do conteúdo, compara-se o comportamento do motor sobre repositórios criados antes e depois do corte de treinamento do modelo base. Dado o volume de repositórios criados na janela, a simulação opera sobre um recorte temático dimensionado pelo orçamento de decisões, produto do número de agentes, do tamanho $L$ da lista de descoberta e do número de passos, de modo a preservar ao menos duas décadas de grau na cauda da distribuição.

O processo de coleta e simulação tem quatro fases. Na Fase I, extração, os eventos são coletados do GH Archive \parencite{grigorik_gharchive} por consulta em plataforma de nuvem, baixados em arquivos Parquet locais e complementados pelos metadados dos repositórios e pelos READMEs versionados, coletados da API oficial do GitHub \parencite{github2025}. Na Fase II, pós-processamento, os READMEs são vetorizados com jina-code-embeddings \parencite{kryvosheieva2025} em banco de dados vetorial indexado por passo mensal, e constrói-se a série temporal de temas de interesse a partir dos tópicos declarados pelos repositórios. Na Fase III, simulação, agentes com personas construídas a partir dos históricos reais de estrelas e fundamentadas no perfil dos desenvolvedores \parencite{borges2018, stackoverflow2025} recebem da plataforma simulada uma lista de descoberta por similaridade semântica de temas e a percorrem segundo o modelo de cascata \parencite{craswell2008}; em seguida, o agente recebe o README completo do resultado escolhido e o motor DS1 devolve a probabilidade de estrela. Cada atribuição de estrela é uma micro-decisão semântica de baixa complexidade e alta frequência, repetida por muitos desenvolvedores sobre muitos repositórios em cada passo temporal, que exige determinismo, custo baixo e rastreabilidade; frente ao LLM generativo da GABM, o DS1 troca julgamento conversacional por classificação de decisão, com menor riqueza contextual, em troca de reprodutibilidade externa das decisões. O motor é o Jebadiah 4B v2 \parencite{frontierinfra2026}, adaptação de pesos abertos sobre o Qwen3.5-4B, executado localmente com quantização Q8 e hash travado. A escolha é justificada por restrição de hardware, pelos pesos abertos, por benchmark técnico independente \parencite{akita2026} e por medição própria em tarefa análoga de classificação temática de artigos científicos, com 784 decisões, 90,3\% de concordância com o banco curado e nenhuma troca de decisão em repetições na mesma máquina. Esse determinismo é condicionado ao ambiente de execução, pois a própria documentação do modelo relata trocas residuais entre repetições e divergências pontuais entre a quantização e os pesos completos \parencite{frontierinfra2026}. A estrela é sorteada por Bernoulli sobre a probabilidade devolvida, com gerador pseudoaleatório da simulação e semente registrada; assim, o motor permanece determinístico, a estocasticidade própria dos modelos de ligação preferencial fica explícita e reprodutível, e as réplicas variam apenas a semente da simulação, com no mínimo 30 sementes por configuração. Para separar o efeito da decisão semântica do efeito da exposição de popularidade, definem-se duas condições: na condição de controle, a lista de descoberta é ordenada apenas por similaridade, com contagem de estrelas oculta; na condição de tratamento, os resultados são reordenados por número de estrelas, como na prática da plataforma. A condição de controle corresponde, assim, a um modelo de aptidão no sentido de \textcite{bianconi2001}, em que a aptidão é o casamento semântico, e a de tratamento combina essa aptidão à exposição por popularidade. Como a ordenação por estrelas e o modelo de cascata favorecem por si sós os repositórios populares, cada condição é executada com três decisores sobre o mesmo processo de descoberta: um decisor aleatório, que concede estrela com a taxa média observada; um decisor por similaridade, que aplica um limiar à similaridade entre os vetores da persona e do README; e o motor DS1. A contribuição do DS1 só se sustenta se seus resultados diferirem de forma mensurável dos dois decisores de referência. Para avaliar a hipótese de custo frente à GABM, uma subamostra das decisões é repetida com o mesmo modelo de base em modo generativo, comparando custo, latência, estabilidade e concordância das decisões. Os parâmetros exógenos, como as taxas de chegada de desenvolvedores e de repositórios e a atividade por desenvolvedor, são calibrados com os dados de 2024, e a validação usa o período de janeiro de 2025 a abril de 2026; a atividade heterogênea dos desenvolvedores entra como dado empírico, pois, sozinha, ela pode produzir cauda pesada.

Na Fase IV, validação, os critérios são hierarquizados. No critério primário, referente ao comportamento de ligação preferencial, registram-se, em cada decisão, o grau do repositório e a lista exibida ao agente. Isso permite decompor a probabilidade de conexão no produto da probabilidade de exposição dado o grau, P(exposto | $k$), pela probabilidade de estrela dada a exposição, P(estrela | exposto, $k$); o segundo fator isola o efeito do motor. Estima-se o expoente $\alpha$ de $\Pi(k) \propto k^{\alpha}$, com intervalo de confiança, pelo método de \textcite{newman2001} e \textcite{jeong2003}, tanto na simulação quanto nos dados do GH Archive, que fornecem o alvo empírico. A equivalência de $\alpha$ a 1 valida a hipótese de que a ligação preferencial linear emerge da decisão semântica; valores entre 0 e 1 ou acima de 1 indicam os regimes sublinear e superlinear; e $\alpha$ próximo de zero refuta-a, indicando emergência por mecanismo distinto. Os critérios secundários reservam o peso da assinatura de escala ao teste estatístico formal, compreendendo o ajuste de Lei da Potência segundo \textcite{clauset2009}, o comparativo com alternativa log-normal \parencite{broido2019}, o Kolmogorov-Smirnov como medida descritiva da distância entre distribuições e o teste de equivalência TOST aplicado aos expoentes $\gamma$ e $\alpha$, com margens registradas em repositório público de pré-registro antes da execução. Em nível microscópico, segue-se a modelagem orientada a padrões \parencite{grimm2005}, pois reproduzir o padrão macroscópico não basta para sustentar as regras geradoras: as personas devem prever as estrelas reais retidas do período de validação, medidas pela taxa de acerto nas primeiras posições, em comparação com os decisores de referência. A análise de sensibilidade varia o número de agentes, o tamanho $L$ da lista de descoberta, o corte de atividade mínima e a redação do \textit{prompt} e dos rótulos das opções \parencite{azizi2026labelsoverridedefinitionsjevstyle}, de modo a garantir que a consistência não dependa de valores pontuais de hiperparâmetros.

Reprodutibilidade é tratada como parte da metodologia, não como anexo, e o protocolo de execução segue o ODD (Visão Geral, Conceitos de Desenho e Detalhes) \parencite{grimm2006}, com versão e hash do modelo travados, ambiente de execução registrado (runtime, hardware, número de \textit{threads} e tamanho de lote), sementes da simulação, \textit{prompts} versionados e logs completos de decisão e de exposição armazenados. Ainda antes da simulação em massa, o motor será submetido a teste de estabilidade, que repete decisões idênticas sob condições idênticas, repete-as com ordem embaralhada dos candidatos e com inversão dos pares apresentados, e mede a taxa de troca e a variação máxima de probabilidade, sendo a sensibilidade à ordem dos candidatos o viés de ordem descrito por \textcite{zhao2021}. Esse teste protege o critério primário, pois, sem ele, a assinatura de ligação preferencial poderia herdar flutuações estocásticas do motor em lugar da decisão semântica. Por fim, a descoberta simulada restringe-se à busca semântica; canais externos e páginas de tendências da plataforma ficam fora do modelo e são tratados como ameaça à validade.

\postextual
\printbibliography[heading=bibintoc,title={Referências Bibliográficas}]
\end{document}
